\section*{Bab \ref{ch:b2:funcseq} --- Barisan dan Deret
Fungsi}

\begin{solution}{exo:b2:funcseq:1}
$f_n(x) = \frac{x}{1 + nx}$: limit \omterm{def:b2:funcseq:def}{titik demi titiknya} $0$ pada
$\intcc{0}{1}$. Keseragamannya: $f_n$ naik pada $\intcc{0}{1}$ (sebab
turunannya $\frac{1}{(1+nx)^2} > 0$), jadi $\norm{f_n}_\infty = f_n(1) =
\frac{1}{1+n} \to 0$: jadi \emph{\omterm{def:b2:funcseq:def}{seragam}} pada $\intcc{0}{1}$.

$g_n(x) = nx\,\eu^{-nx^2}$: limit \omterm{def:b2:funcseq:def}{titik demi titiknya} $0$ (sebab
eksponensialnya menang). Supremumnya: $g_n' = n\eu^{-nx^2}(1 - 2nx^2)$
nol di $x_n = \frac{1}{\sqrt{2n}}$, tempat $g_n(x_n) =
\sqrt{\frac n2}\,\eu^{-1/2} \to \infty$: jadi tidak \omterm{def:b2:funcseq:def}{seragam} pada
$\intcc{0}{1}$ --- tetapi \omterm{def:b2:funcseq:def}{seragam} pada $\intco{\delta}{1}$, sebab di
sana $g_n(x) \leq n\,\eu^{-n\delta^2} \to 0$.

$h_n(x) = x^n(1 - x^n)$: limit \omterm{def:b2:funcseq:def}{titik demi titiknya} $0$ pada
$\intcc{0}{1}$ (dari kedua faktornya; dan di $x = 1$, $h_n = 0$).
Supremumnya: dengan $u = x^n \in \intcc{0}{1}$, berlaku $u(1-u) \leq
\frac14$ yang tercapai di $u = \frac12$, yakni $x = 2^{-1/n} \in
\intoo{0}{1}$: jadi $\norm{h_n}_\infty = \frac14 \not\to 0$, sehingga
tidak \omterm{def:b2:funcseq:def}{seragam} pada $\intcc{0}{1}$; tetapi \omterm{def:b2:funcseq:def}{seragam} pada
$\intcc{0}{a}$ (sebab $\sup \leq a^n \to 0$).
\end{solution}

\begin{solution}{exo:b2:funcseq:2}
Berlaku $\int_0^1 nx\,\eu^{-nx^2}\dd x = \bigl[-\tfrac12
\eu^{-nx^2}\bigr]_0^1 = \frac{1 - \eu^{-n}}{2} \to \frac12$, padahal
$\int_0^1 \lim g_n = 0$. Tak ada kontradiksi:
\cref{thm:b2:funcseq:integration} menuntut kekonvergenan \emph{\omterm{def:b2:funcseq:def}{seragam}}
pada ruas itu, dan di sini hal itu gagal (sebab bonggol setinggi
$\sim\sqrt n$ meluncur ke $0$).
\end{solution}

\begin{solution}{exo:b2:funcseq:3}
\omterm{def:b2:funcseq:series}{Kekonvergenan normal} pada $\intcc{-1}{1}$: $\norm{x^n/n^2}_\infty =
\frac{1}{n^2}$, yang \omterm{def:b2:series:summable}{terjumlahkan}: jadi $S$ \omterm{def:b2:metric:continuity}{kontinu} di sana
(\cref{thm:b2:funcseq:seriestransfer}).

Turunannya: deret turunannya $\sum \frac{x^{n-1}}{n}$ konvergen
\omterm{def:b2:funcseq:series}{normal} pada tiap $\intcc{-a}{a}$ dengan $a < 1$ (sebab $\sup =
\frac{a^{n-1}}{n}$): jadi $S$ berkelas $C^1$ pada $\intoo{-1}{1}$ dengan
\[
S'(x) = \sum_{n\geq1} \frac{x^{n-1}}{n} = \frac1x \sum_{n\geq1}
\frac{x^n}{n} = -\frac{\ln(1 - x)}{x}
\qquad (0 < \abs x < 1),
\]
dengan kesamaan terakhirnya adalah deret logaritmik Tahun ke-1 (yang
diturunkan secara jujur pada \cref{ch:b2:powerseries}).
\end{solution}

\begin{solution}{exo:b2:funcseq:4}
Untuk $x > 0$ yang tetap, deretnya berselang-seling dengan $\frac{1}{n +
x} \downarrow 0$: jadi konvergen \omterm{def:b2:funcseq:def}{titik demi titik}, dan batas sisanya
$\abs{R_N(x)} \leq \frac{1}{N + 1 + x} \leq \frac{1}{N+1}$ bersifat
\emph{\omterm{def:b2:funcseq:def}{seragam}} pada $\intco{\delta}{\infty}$ (bahkan pada
$\intoo{0}{\infty}$): jadi kekonvergenannya \omterm{def:b2:funcseq:def}{seragam}. (Tidak \omterm{def:b2:funcseq:series}{normal}:
sebab $\norm{\frac{(-1)^n}{n+x}}_\infty = \frac{1}{n + \delta}$, yang
divergen.) Kekontinuannya menyusul dari
\cref{thm:b2:funcseq:seriestransfer}.

Persamaan fungsionalnya: indeks ulang $F(x + 1)$ dengan $m = n + 1$:
\[
F(x+1) = \sum_{n\geq0} \frac{(-1)^n}{n + 1 + x}
= \sum_{m\geq1}\frac{(-1)^{m-1}}{m+x} ,
\]
jadi, dengan memisahkan suku $m = 0$ pada $F(x)$,
\[
F(x) + F(x+1)
= \frac{1}{x} + \sum_{m\geq1}
\frac{(-1)^m + (-1)^{m-1}}{m+x} = \frac1x .
\]
\end{solution}

\begin{solution}{exo:b2:funcseq:5}
Tetapkan $g_n = f_n - f \geq 0$ (yang turun terhadap $n$ menurut
hipotesisnya; dan limitnya $0$ \omterm{def:b2:funcseq:def}{titik demi titik}); tiap $g_n$ \omterm{def:b2:metric:continuity}{kontinu}.
Tetapkan $\varepsilon > 0$ lalu misalkan $U_n = \{x : g_n(x) <
\varepsilon\}$: ia \omterm{def:b2:metric:topology}{terbuka} (sebab prapeta himpunan \omterm{def:b2:metric:topology}{terbuka}), naik (sebab
$g_{n+1} \leq g_n$), dan menyelimuti $K$ (berkat \omterm{def:b2:funcseq:def}{kekonvergenan titik
demi titik}). Menurut Borel--Lebesgue
(\cref{thm:b2:metric:borellebesgue}), berhingga banyak $U_{n_1}
\subseteq \dots \subseteq U_{n_k}$ menyelimuti $K$: jadi $K = U_{n_k}$,
yakni $\norm{g_{n_k}}_\infty \leq \varepsilon$, dan berkat kemonotonannya
$\norm{g_n}_\infty \leq \varepsilon$ untuk setiap $n \geq n_k$: jadi
kekonvergenannya \omterm{def:b2:funcseq:def}{seragam}. (Kemonotonannya penting: sebab bonggol
meluncur pada \cref{exo:b2:funcseq:2} konvergen \omterm{def:b2:funcseq:def}{titik demi titik} pada
himpunan \omterm{def:b2:metric:compact}{kompak} tanpa keseragaman.)
\end{solution}

\begin{solution}{exo:b2:funcseq:6}
Sulihkan $u = nx$:
\[
\int_0^1 \frac{n f(x)}{1 + n^2x^2}\dd x
= \int_0^n \frac{f(u/n)}{1 + u^2}\,\dd u .
\]
Integran $h_n(u) = \frac{f(u/n)}{1+u^2}\mathbf{1}_{u \leq n}$
konvergen \omterm{def:b2:funcseq:def}{titik demi titik} ke $\frac{f(0)}{1+u^2}$ (berkat kekontinuan
$f$ di $0$) dan terdominasi oleh $\frac{\norm f_\infty}{1 + u^2}$, yang
terintegralkan pada $\intco{0}{\infty}$: jadi kekonvergenan terdominasi
(\cref{thm:b2:integration:dominated}) memberi limitnya
\[
\int_0^\infty \frac{f(0)}{1 + u^2}\dd u = \frac{\pi}{2} f(0) .
\]
(Kernelnya memusat di $0$: yakni sebuah identitas hampiran.)
\end{solution}

\begin{solution}{exo:b2:funcseq:7}
Untuk $f(x) = x^2$, pakailah keluarga kesamaan binomial kedua dari
buktinya: $\sum_k k^2 p_k = n(n-1)x^2 + nx$. Karenanya
\[
B_n(f)(x) = \sum_k \frac{k^2}{n^2}\,p_k
= \frac{n(n-1)x^2 + nx}{n^2}
= x^2 + \frac{x(1 - x)}{n} :
\]
jadi $B_n(f) - f = \frac{x(1-x)}{n}$, yang bernorma supremum
$\frac{1}{4n} = O\bigl(\frac1n\bigr)$, sesuai ramalannya.
\end{solution}

\begin{solution}{exo:b2:funcseq:8}
Untuk $\varepsilon = 1$ ada $N$ dengan $\norm{P_n - P_m}_{\infty,
\R} \leq 1$ bagi $m, n \geq N$. Polinomial yang terbatas pada $\R$
pastilah konstan (sebab yang tak konstan menuju $\pm\infty$): jadi $P_n
- P_m = c_{n,m}$, yakni konstanta. Karenanya untuk $n \geq N$: $P_n =
P_N + c_n$ dengan $c_n = P_n(0) - P_N(0)$ yang konvergen (berkat
\omterm{def:b2:funcseq:def}{kekonvergenan titik demi titik} di $0$). Karenanya $f = \lim P_n = P_N +
\lim c_n$: sebuah polinomial.
\end{solution}

\begin{solution}{exo:b2:funcseq:9}
(a) Di sini $\norm{(3/4)^n\varphi(4^n\cdot)}_\infty = \frac12
(3/4)^n$: jadi kekonvergenannya \omterm{def:b2:funcseq:series}{normal}, sehingga $W$ \omterm{def:b2:metric:continuity}{kontinu}
(\cref{thm:b2:funcseq:seriestransfer}).

(b) Tetapkan $x$ dan $m$; pilihlah tanda $h_m = \pm\frac12 4^{-m}$
sehingga ruas $\intcc{4^mx}{4^m(x + h_m)}$ (yang berpanjang $\frac12$)
tak memuat satu pun bilangan setengah bulat, sehingga $\varphi$ menjadi
afin berkemiringan $\pm1$ di sana (hal ini mungkin: sebab interval
berpanjang $\frac12$ bertemu paling banyak satu titik setengah bulat;
pilihlah sisi yang menghindarinya).

Untuk $n > m$: $4^n h_m = \pm\frac12 4^{n-m}$ merupakan bilangan bulat,
dan $\varphi$ berperiode $1$: jadi suku ke-$n$ pada selisihnya lenyap.

Untuk $n = m$: $\abs{\varphi(4^m x + 4^m h_m) - \varphi(4^m x)} =
\abs{4^m h_m} = \frac12$ (sebab $\varphi$ afin berkemiringan $\pm 1$ pada
ruas itu), jadi sukunya menyumbang tepat $(3/4)^m \cdot
\frac{1/2}{\abs{h_m}} = (3/4)^m\,4^m = 3^m$ pada hasil baginya.

Untuk $n < m$: sifat Lipschitz-$1$ pada $\varphi$ memberi
$\bigl|(3/4)^n\bigl(\varphi(4^nx + 4^nh_m) -
\varphi(4^nx)\bigr)\bigr| \leq (3/4)^n 4^n\abs{h_m} =
3^n\abs{h_m}$: jadi masing-masing menyumbang paling banyak $3^n$ pada
hasil baginya.

Karenanya
\[
\Bigl|\frac{W(x + h_m) - W(x)}{h_m}\Bigr|
\geq 3^m - \sum_{n=0}^{m-1} 3^n
= 3^m - \frac{3^m - 1}{2} = \frac{3^m + 1}{2}
\longrightarrow \infty .
\]
Seandainya $W$ dapat diturunkan di $x$, maka setiap hasil bagi selisih
sepanjang $h_m \to 0$ akan konvergen ke $W'(x)$: kontradiksi. Jadi $W$
\omterm{def:b2:metric:continuity}{kontinu} di mana-mana dan tak dapat diturunkan di mana pun.
\end{solution}

\begin{solution}{exo:b2:funcseq:10}
\omterm{def:b2:funcseq:def}{Titik demi titik}: untuk $x \in \intco{0}{1}$ deretnya geometri dengan
rasio $-x$,
\[
\sum_{n\geq0}(-1)^n x^n(1-x) = \frac{1-x}{1+x},
\]
sedangkan di $x = 1$ tiap sukunya nol, jadi jumlahnya $0 =
\frac{1-1}{2}$, yang selaras --- sehingga jumlahnya \omterm{def:b2:metric:continuity}{kontinu} pada
$\intcc{0}{1}$. Keseragamannya: sisanya berupa ekor geometri,
\[
\abs{R_N(x)} = \frac{x^{N+1}(1-x)}{1+x} \leq x^{N+1}(1 - x)
\leq \max_{\intcc01} t^{N+1}(1-t)
= \frac{1}{N+2}\Bigl(\frac{N+1}{N+2}\Bigr)^{\!N+1}
\leq \frac{1}{N+2} \to 0 ,
\]
secara \omterm{def:b2:funcseq:def}{seragam} pada $x$. Tidak \omterm{def:b2:funcseq:series}{normal}: sebab $\norm{u_n}_\infty =
\max x^n(1-x) = \frac{1}{n+1}\bigl(\frac{n}{n+1}\bigr)^n \sim
\frac{1}{\eu\,n}$, dan $\sum \frac1{\eu n}$ divergen. Bandingkan:
$\sum x^n(1-x)$ punya jumlah parsial $1 - x^{N+1}$, yang konvergen
\omterm{def:b2:funcseq:def}{titik demi titik} ke $\mathbf 1_{\intco01}$ yang \emph{tak \omterm{def:b2:metric:continuity}{kontinu}}: jadi
menurut \cref{thm:b2:funcseq:continuity}, kekonvergenan itu tak mungkin
seragam pada $\intcc{0}{1}$.
\end{solution}

\begin{solution}{exo:b2:funcseq:11}
Limitnya $f$ bersifat \omterm{def:b2:metric:continuity}{kontinu} (\cref{thm:b2:funcseq:continuity}).
Maka
\[
\abs{f_n(x_n) - f(x)}
\leq \abs{f_n(x_n) - f(x_n)} + \abs{f(x_n) - f(x)}
\leq \norm{f_n - f}_\infty + \abs{f(x_n) - f(x)} ,
\]
dan kedua sukunya menuju $0$ (berkat \omterm{def:b2:funcseq:def}{kekonvergenan seragam} dan
kekontinuan $f$ di $x$). Contoh penyangkal di bawah \omterm{def:b2:funcseq:def}{kekonvergenan titik
demi titik} semata: $g_n(x) = nx\,\eu^{-nx^2} \to 0$ \omterm{def:b2:funcseq:def}{titik demi titik} pada
$\intcc{0}{1}$ dengan $g_n$ dan limitnya \omterm{def:b2:metric:continuity}{kontinu}, namun di $x_n =
\frac{1}{\sqrt{2n}} \to 0$:
\[
g_n(x_n) = \sqrt{\frac n2}\,\eu^{-1/2} \longrightarrow +\infty
\neq 0 = f(0) .
\]
\end{solution}

\begin{solution}{exo:b2:funcseq:12}
\begin{enumerate}
  \item Induksi. Kasus $n = 1$ adalah definisinya. Andaikan rumusnya
        berlaku untuk $n$ lalu tetapkan $g(x) = \int_0^x
        \frac{(x-t)^n}{n!} f(t)\dd t$. Untuk integran yang \omterm{def:b2:metric:continuity}{kontinu} pada
        $(x,t)$ dan berkelas $C^1$ pada $x$, \omterm{thm:b2:integration:continuity}{integral berparameter}
        dengan batas peubah itu terturunkan sebagai
        \[
        g'(x) = \frac{(x-x)^n}{n!}f(x)
        + \int_0^x \frac{(x-t)^{n-1}}{(n-1)!}f(t)\dd t
        = T^nf(x)
        \]
        (pecahlah $g(x+h) - g(x)$ menjadi jalur
        $\int_x^{x+h}$, yang bernilai $O(h\cdot\sup)$ dengan
        integrannya nol di $t = x$ seperti $h^n$, dan integral tetap
        atas pertambahan $x$-nya, yang ditangani lewat ketaksamaan nilai
        rata-rata dan kekontinuan). Selain itu $(T^{n+1}f)' = T^nf$
        (lewat teorema dasar kalkulus) dan $g(0) = T^{n+1}f(0) = 0$: jadi
        dua antiturunan $T^nf$ yang nol di $0$ berimpit, sehingga
        $T^{n+1}f = g$. Adapun batasnya:
        \[
        \abs{T^nf(x)} \leq \norm f_\infty
        \int_0^x \frac{(x-t)^{n-1}}{(n-1)!}\dd t
        = \norm f_\infty\,\frac{x^n}{n!}
        \leq \frac{\norm f_\infty}{n!} .
        \]
  \item Berlaku $\sum_n \norm{T^nf}_\infty \leq \eu\,\norm f_\infty$:
        jadi kekonvergenannya \omterm{def:b2:funcseq:series}{normal}, sehingga \omterm{def:b2:funcseq:def}{seragam}; dan $S$ \omterm{def:b2:metric:continuity}{kontinu}.
        Jumlah parsialnya memenuhi $S_N = f + T S_{N-1}$, sedangkan $T$
        bersifat Lipschitz-$1$ terhadap $\norm\cdot_\infty$ (sebab
        $\abs{Tg(x)} \leq x\norm g_\infty$): jadi melewatkan $N \to
        \infty$ pada kedua ruasnya memberi $S = f + TS$.
  \item Tetapkan $V(x) = f(x) + \eu^x\int_0^x \eu^{-t}f(t)\dd t$.
        Maka $V - f$ berkelas $C^1$ dengan $(V-f)'(x) =
        \eu^x\int_0^x\eu^{-t}f + f(x) = V(x)$, sedangkan $(TV)' = V$
        dengan $(V - f)(0) = TV(0) = 0$: jadi $V - f = TV$, yakni
        $V$ menyelesaikan persamaannya. Ketunggalannya: jika $S_1, S_2$
        dua penyelesaian yang \omterm{def:b2:metric:continuity}{kontinu}, maka $D = S_1 - S_2$ memenuhi $D
        = TD$, sehingga $D = T^nD$ untuk setiap $n$ dan $\norm D_\infty
        \leq \frac{\norm D_\infty}{n!} \to 0$: jadi $D = 0$.
        Karenanya
        \[
        \sum_{n\geq0} T^nf(x) = f(x) +
        \int_0^x \eu^{x-t}f(t)\,\dd t .
        \]
        (Deret $\sum T^n$ adalah deret geometri atas
        operator: yakni cicipan pertama resolven
        $(\mathrm{Id} - T)^{-1}$, yang dikembangkan pada jilid Tahun
        ke-3.)
\end{enumerate}
\end{solution}

\begin{solution}{pb:b2:funcseq:1}
\textbf{1.} Kelinearannya jelas dari rumusnya. Kepositifannya:
bobot $p_k(x) \geq 0$, jadi $f \geq 0$ memaksa $B_nf \geq
0$; dan kemonotonannya menyusul setelah diterapkan pada $g - f$.
Batasnya: dari $\pm f \leq \norm f_\infty$ diperoleh $\pm B_nf \leq
\norm f_\infty B_ne_0 = \norm f_\infty$. Di ujungnya: $p_k(0) =
\mathbf 1_{k=0}$ dan $p_k(1) = \mathbf 1_{k=n}$, jadi $B_nf(0) = f(0)$
dan $B_nf(1) = f(1)$.

\textbf{2.} Turunkan $(x+y)^n = \sum_k\binom nk x^ky^{n-k}$
terhadap $x$, kalikan dengan $x$, lalu ambil $y = 1 - x$:
\[
nx = \sum_k k\,p_k(x) ;
\]
dua kali, lalu dikalikan $x^2$: $n(n-1)x^2 = \sum_k k(k-1)p_k(x)$.
Karenanya $B_ne_0 = 1$ (lewat teorema binomial), $B_ne_1(x) =
\frac{nx}{n} = x$, dan
\[
B_ne_2(x) = \frac{\sum_k k^2p_k}{n^2}
= \frac{n(n-1)x^2 + nx}{n^2}
= x^2 + \frac{x(1-x)}{n} .
\]

\textbf{3.} Uraikan:
\[
\sum_k\Bigl(\frac kn - x\Bigr)^{\!2} p_k
= B_ne_2(x) - 2x\,B_ne_1(x) + x^2
= \frac{x(1-x)}{n} .
\]
Lalu Cauchy--Schwarz dengan pemecahan $\abs{k/n - x}\sqrt{p_k}
\cdot \sqrt{p_k}$ memberi
\[
\sum_k\Bigl|\frac kn - x\Bigr| p_k
\leq \Bigl(\sum_k\Bigl(\frac kn - x\Bigr)^2
p_k\Bigr)^{\!1/2}
= \sqrt{\frac{x(1-x)}{n}} \leq \frac{1}{2\sqrt n},
\]
dengan memakai $x(1-x) \leq \frac14$.

\textbf{4.} Bobot $p_k(x)$ bersifat tak negatif dengan jumlah $1$
dan barisenter $\sum_k \frac kn p_k(x) = x$ (pertanyaan 2). Lalu
ketaksamaan Jensen yang hingga bagi $f$ yang cembung (lewat induksi dari
definisi dua titiknya, jilid Tahun ke-1) memberi
\[
f(x) = f\Bigl(\sum_k \frac kn\,p_k\Bigr)
\leq \sum_k f\Bigl(\frac kn\Bigr)p_k = B_nf(x) .
\]

\textbf{5.} Pada $\{k : \abs{k/n - x} > \delta\}$ berlaku
$\bigl(\frac{k/n - x}{\delta}\bigr)^2 > 1$, jadi
\[
\sum_{\abs{k/n-x}>\delta} p_k
\leq \frac{1}{\delta^2}\sum_k\Bigl(\frac kn -
x\Bigr)^2p_k
= \frac{x(1-x)}{n\delta^2} \leq \frac{1}{4n\delta^2} .
\]
Pembacaannya: $p_k(x)$ adalah distribusi frekuensi sampel $S_n/n$ atas
$n$ lemparan koin berbias $x$; nilai harapannya $x$, ragamnya
$\frac{x(1-x)}n \to 0$, dan ungkapan di atas tak lain ketaksamaan
Chebyshev: jadi massanya memusat di $x$, sehingga merata-ratakan $f$
terhadapnya memproduksi kembali $f(x)$ pada limitnya
(\cref{ch:b2:genfun} meresmikan kosakatanya).

\textbf{6.} Berlaku $\omega \leq 2M < \infty$; kemonotonannya jelas
(sebab supremum atas himpunan yang lebih besar). Heine: $f$ yang \omterm{def:b2:metric:continuity}{kontinu}
pada himpunan \omterm{def:b2:metric:compact}{kompak} bersifat kontinu seragam, dan itu persis berbunyi
$\omega(\delta) \to 0$ ketika $\delta \to 0^+$. Subaditifnya: jika
$\abs{s - t} \leq \delta_1 + \delta_2$, maka titik $u$ pada ruas
$\intcc st$ yang berjarak $\min(\delta_1, \abs{s-t})$ dari $s$ memenuhi
$\abs{s-u} \leq \delta_1$ dan $\abs{u-t} \leq \delta_2$, sedangkan
$\abs{f(s)-f(t)} \leq \abs{f(s)-f(u)} + \abs{f(u)-f(t)}$.
Setelah diiterasikan, $\omega(p\delta) \leq p\,\omega(\delta)$ untuk $p
\in \N^*$; lalu untuk $\lambda > 0$, dengan $p = \lceil\lambda\rceil
\leq 1 + \lambda$: $\omega(\lambda\delta) \leq \omega(p\delta) \leq
p\,\omega(\delta) \leq (1+\lambda)\omega(\delta)$.

\textbf{7.} Karena $\sum p_k = 1$:
\[
\abs{B_nf(x) - f(x)}
= \Bigl|\sum_k\bigl(f(k/n) - f(x)\bigr)p_k\Bigr|
\leq \sum_k\omega\bigl(\abs{k/n - x}\bigr)p_k .
\]
Untuk tiap $k$, pertanyaan 6 dengan $\lambda = \abs{k/n - x}/\delta$
memberi $\omega(\abs{k/n-x}) \leq \bigl(1 +
\frac{\abs{k/n-x}}\delta\bigr)\omega(\delta)$; lalu menjumlahkannya
terhadap $p_k$ menghasilkan taksiran induknya.

\textbf{8.} Masukkan batas pertanyaan 3:
\[
\abs{B_nf(x) - f(x)} \leq \Bigl(1 +
\frac{1}{2\delta\sqrt n}\Bigr)\omega(\delta),
\]
secara \omterm{def:b2:funcseq:def}{seragam} pada $x$; lalu dengan $\delta = n^{-1/2}$ kurungnya
menjadi $\frac32$: jadi $\norm{B_nf - f}_\infty \leq
\frac32\omega(n^{-1/2}) \to 0$ menurut pertanyaan 6 (Heine). Ini
membuktikan ulang \cref{thm:b2:funcseq:weierstrass} beserta lajunya.

\textbf{9.} Lipschitz-$L$ berarti $\omega(\delta) \leq L\delta$:
jadi lajunya $\frac{3L}{2\sqrt n}$. Sedangkan Hölder-$\alpha$ berarti
$\omega(\delta) \leq C\delta^\alpha$: jadi lajunya $\frac{3C}{2}
n^{-\alpha/2}$.

\textbf{10.} Dengan memakai $k\binom{2m}k = 2m\binom{2m-1}{k-1}$:
\[
\sum_{k=m+1}^{2m}k\binom{2m}k
= 2m\sum_{j=m}^{2m-1}\binom{2m-1}{j}
= 2m\cdot 2^{2m-2},
\]
sebab $j \mapsto 2m-1-j$ membijeksikan $\{m,\dots,2m-1\}$ pada
$\{0,\dots,m-1\}$, sehingga jumlahnya separuh dari $2^{2m-1}$. Selain
itu $\sum_{k=m+1}^{2m}\binom{2m}k = \frac{2^{2m} -
\binom{2m}m}{2}$ (lewat kesetangkupan yang sama). Karenanya
\[
\sum_{k=m+1}^{2m}(k-m)\binom{2m}k
= m\,2^{2m-1} - m\,\frac{2^{2m} - \binom{2m}m}{2}
= \frac m2\binom{2m}m .
\]
Penyulihan $k \mapsto 2m-k$ memetakan suku dengan $k < m$
pada suku dengan $k > m$ (dengan binomial yang sama dan $\abs{k-m}$ yang
sama): jadi jumlah \omterm{def:b2:series:def}{mutlaknya} dua kali jumlah sepihaknya, yakni $m\binom{2m}m$.

\textbf{11.} Di $x = \frac12$ berlaku $p_k(\tfrac12) =
\binom{2m}k2^{-2m}$ dan $f(\tfrac12) = 0$, jadi
\[
B_{2m}f\Bigl(\frac12\Bigr)
= \sum_k\Bigl|\frac{k}{2m} - \frac12\Bigr|
\binom{2m}k 2^{-2m}
= \frac{2^{-2m}}{2m}\,m\binom{2m}m
= \frac{\binom{2m}m}{2\cdot4^m}
\sim \frac{1}{2\sqrt{\pi m}}
\]
menurut \cref{ex:b2:comparison:centralbinomial}. Karena $\omega_f
(\delta) = \delta$ di sini (sebab fungsinya Lipschitz-$1$ dan
batasnya tercapai), pertanyaan 8 meramalkan paling banyak
$\frac32(2m)^{-1/2}$: sedangkan galat sejatinya $\frac{1}{2\sqrt{\pi m}}$
tepat berorde $n^{-1/2}$ --- jadi lajunya tajam kecuali
konstantanya.

\textbf{12.} Dari $f \leq g$ diperoleh $g - f \geq 0$, jadi $L(g-f) \geq
0$, yakni $Lf \leq Lg$. Dari $-\abs f \leq f \leq \abs f$:
$-L\abs f \leq Lf \leq L\abs f$, yakni $\abs{Lf} \leq
L\abs f$.

\textbf{13.} Lewat Heine pilihlah $\delta$ dengan $\abs{f(s)-f(x)}
\leq \varepsilon$ setiap kali $\abs{s-x} \leq \delta$. Jika
$\abs{s - x} > \delta$, maka $\frac{(s-x)^2}{\delta^2} > 1$ dan
$\abs{f(s)-f(x)} \leq 2M \leq \frac{2M}{\delta^2}(s-x)^2$. Jadi pada
kedua kasusnya batas yang diklaim itu berlaku.

\textbf{14.} Tetapkan $x$; pertanyaan 13 berbunyi, sebagai fungsi $s$:
\[
-\varepsilon e_0 - \frac{2M}{\delta^2}q_x
\;\leq\; f - f(x)e_0
\;\leq\; \varepsilon e_0 + \frac{2M}{\delta^2}q_x,
\qquad q_x = e_2 - 2x\,e_1 + x^2e_0 .
\]
Lalu terapkan $L_n$ yang linear monoton (pertanyaan 12) lalu nilailah di
$x$:
\[
\abs{L_nf(x) - f(x)L_ne_0(x)}
\leq \varepsilon L_ne_0(x)
+ \frac{2M}{\delta^2}\bigl(L_ne_2(x) - 2xL_ne_1(x)
+ x^2L_ne_0(x)\bigr) .
\]

\textbf{15.} Tulis $\alpha_j = L_ne_j - e_j$, sehingga
$\norm{\alpha_j}_\infty \to 0$. Karena $e_2(x) - 2xe_1(x) +
x^2e_0(x) = 0$:
\[
L_ne_2(x) - 2xL_ne_1(x) + x^2L_ne_0(x)
= \alpha_2(x) - 2x\,\alpha_1(x) + x^2\alpha_0(x),
\]
yang bernorma supremum paling banyak $\norm{\alpha_2} +
2\norm{\alpha_1} + \norm{\alpha_0} \to 0$. Selain itu $L_ne_0 \to e_0$
secara \omterm{def:b2:funcseq:def}{seragam}, jadi $L_ne_0 \leq 2$ untuk $n$ yang besar, dan
$\abs{f(x)}\abs{L_ne_0(x) - 1} \leq M\norm{\alpha_0} \to 0$. Setelah
dirakit dengan pertanyaan 14: untuk $n$ yang besar, \omterm{def:b2:funcseq:def}{seragam} pada $x$,
\[
\abs{L_nf(x) - f(x)} \leq 2\varepsilon +
\frac{2M}{\delta^2}\,o(1) + M\,o(1) \leq 3\varepsilon :
\]
jadi $L_nf \to f$ secara \omterm{def:b2:funcseq:def}{seragam} --- itulah teorema Korovkin.

\textbf{16.} Berlaku $B_ne_0 = e_0$ dan $B_ne_1 = e_1$ secara eksak, dan
$\norm{B_ne_2 - e_2}_\infty = \max_x\frac{x(1-x)}{n} =
\frac{1}{4n} \to 0$ (pertanyaan 2): jadi Korovkin berlaku, dan
Weierstrass menyusul untuk ketiga kalinya.

\textbf{17.} Pemetaan $I_nf$ linear pada $f$ (sebab nilai simpulnya
demikian), dan pada tiap sel penginterpolasi afin atas nilai simpul yang
tak negatif bersifat tak negatif: jadi positif. Lalu $I_ne_0 = e_0$ dan
$I_ne_1 = e_1$ sebab fungsi afin sama dengan penginterpolasinya sendiri.
Pada sebuah sel $\intcc ab$ (dengan $b - a = \frac1n$), penginterpolasi
afin bagi $e_2$ adalah $L(t) = (a+b)t - ab$, dan
\[
L(t) - t^2 = (t-a)(b-t) \in
\intcc{0}{\tfrac{(b-a)^2}{4}} ,
\]
dengan maksimumnya di titik tengahnya: jadi $\norm{I_ne_2 - e_2}_\infty =
\frac{1}{4n^2} \to 0$. Menurut Korovkin: $I_nf \to f$ secara \omterm{def:b2:funcseq:def}{seragam} untuk
setiap $f$ yang \omterm{def:b2:metric:continuity}{kontinu} --- yakni hampiran poligonal, tanpa perlu
taksiran lebih lanjut.

\textbf{18.} Diberikan $f$ dan $\varepsilon$: Weierstrass menyediakan
polinomial $P = \sum_{j=0}^d a_jx^j$ dengan $\norm{f - P}_\infty \leq
\frac\varepsilon2$; lalu mengganti tiap $a_j$ dengan bilangan rasional
$b_j$ yang $\abs{a_j - b_j} \leq \frac{\varepsilon}{2(d+1)}$ menggeser
\omterm{def:b2:nvs:norm}{norma} supremumnya pada $\intcc{0}{1}$ paling banyak
$\frac\varepsilon2$. Adapun himpunan polinomial berkoefisien rasional
adalah gabungan \omterm{def:b2:structures:countable}{terbilang} (atas $d$) atas \omterm{def:b2:structures:countable}{himpunan terbilang}, jadi
\omterm{def:b2:structures:countable}{terbilang}, dan sekaligus padat: sehingga $C(\intcc{0}{1})$ terpisahkan.

\textbf{19.} Berkat kelinearan, $\int_0^1 fP = 0$ untuk setiap
polinomial $P$. Pilihlah polinomial $P_n \to f$ yang konvergen \omterm{def:b2:funcseq:def}{seragam}
(lewat Weierstrass):
\[
\Bigl|\int_0^1 f^2\Bigr|
= \Bigl|\int_0^1 f\,(f - P_n)\Bigr|
\leq \norm f_\infty\,\norm{f - P_n}_\infty
\longrightarrow 0 ,
\]
jadi $\int_0^1 f^2 = 0$. Seandainya $f(x_0) \neq 0$, kekontinuan memberi
$f^2 \geq c > 0$ pada sebuah subinterval, yang bertentangan dengan
penolkan integralnya: jadi $f = 0$. Akibatnya dua fungsi \omterm{def:b2:metric:continuity}{kontinu} dengan
momen $\int f t^n$ yang sama pastilah berimpit.

\textbf{20.} Dengan $p_{n,k}(x) = \binom nk x^k(1-x)^{n-k}$ dan
kesepakatan $p_{n-1,-1} = p_{n-1,n} = 0$, aturan hasil kali
beserta $k\binom nk = n\binom{n-1}{k-1}$ dan $(n-k)\binom nk =
n\binom{n-1}{k}$ memberi
\[
p_{n,k}'(x) = n\bigl(p_{n-1,k-1}(x) - p_{n-1,k}(x)\bigr) .
\]
Lalu menjumlahkannya terhadap $f(k/n)$ dan menggeser indeks pada jumlah
pertamanya (yakni \omterm{thm:b2:series:abel}{penjumlahan Abel}) memberi
\[
(B_nf)'(x) = n\sum_{j=0}^{n-1}\Bigl(f\Bigl(\frac{j+1}n\Bigr)
- f\Bigl(\frac jn\Bigr)\Bigr)p_{n-1,j}(x) .
\]

\textbf{21.} Menurut teorema nilai rata-rata, $f(\frac{j+1}n) -
f(\frac jn) = \frac1n f'(\xi_j)$ dengan $\xi_j \in
\intoo{j/n}{(j+1)/n}$, jadi $(B_nf)'(x) = \sum_j
f'(\xi_j)\,p_{n-1,j}(x)$. Simpul $\frac{j}{n-1}$ juga terletak di
$\intcc{j/n}{(j+1)/n}$ (sebab kedua ketaksamaannya menyusut menjadi $j
\leq n-1$), sehingga $\abs{\xi_j - \frac j{n-1}} \leq \frac1n$ dan
\[
\bigl|(B_nf)'(x) - B_{n-1}(f')(x)\bigr|
\leq \sum_j\Bigl|f'(\xi_j) -
f'\Bigl(\frac{j}{n-1}\Bigr)\Bigr| p_{n-1,j}(x)
\leq \omega_{f'}\Bigl(\frac1n\Bigr) \longrightarrow 0
\]
secara \omterm{def:b2:funcseq:def}{seragam}. Karena $B_{n-1}(f') \to f'$ secara \omterm{def:b2:funcseq:def}{seragam}
(\cref{thm:b2:funcseq:weierstrass} yang diterapkan pada $f'$ yang
\omterm{def:b2:metric:continuity}{kontinu}), ketaksamaan segitiga memberi $(B_nf)' \to f'$ secara
\omterm{def:b2:funcseq:def}{seragam}. Karenanya polinomial $P_n = B_nf$ konvergen ke $f$ dalam
arti $C^1$.

\textbf{22.} Taylor--Lagrange di $x$: $f(\frac kn) - f(x) =
f'(x)(\frac kn - x) + \frac{f''(\xi_k)}2(\frac kn - x)^2$.
Setelah dijumlahkan terhadap $p_k$, suku linearnya mati (pertanyaan 2):
\[
\abs{B_nf(x) - f(x)}
\leq \frac{\norm{f''}_\infty}{2}\sum_k\Bigl(\frac kn -
x\Bigr)^2p_k
= \frac{\norm{f''}_\infty}{2}\cdot\frac{x(1-x)}{n}
\leq \frac{\norm{f''}_\infty}{8n} .
\]

\textbf{23.} Dua kali penurunan lagi atas $(x+y)^n$ memberi
momen faktorialnya, dengan $n_{(j)} = n(n-1)\cdots(n-j+1)$:
\[
\sum_k k_{(j)}\,p_k = n_{(j)}\,x^j \qquad (j = 3, 4),
\]
sedangkan $k^3 = k_{(3)} + 3k_{(2)} + k$ dan $k^4 = k_{(4)} + 6k_{(3)} +
7k_{(2)} + k$ mengubahnya menjadi momen pangkat:
\[
\sum_k k^3p_k = n_{(3)}x^3 + 3n_{(2)}x^2 + nx,
\qquad
\sum_k k^4p_k = n_{(4)}x^4 + 6n_{(3)}x^3 + 7n_{(2)}x^2 + nx .
\]
Setelah $(k - nx)^4$ diuraikan lalu dikumpulkan (yakni perhitungan yang
sabar tetapi murni mekanis dengan keempat momen pangkat tadi):
\[
\sum_k(k-nx)^4p_k = nx(1-x)\bigl(1 + 3(n-2)x(1-x)\bigr) .
\]
Dengan $x(1-x) \leq \frac14$: ruas kanannya paling banyak
$\frac n4\bigl(1 + \frac{3n}4\bigr) = \frac{3n^2}{16} +
\frac n4 \leq n^2$ untuk $n \geq 1$.

\textbf{24.} Bentuk Peano pada Taylor di $x$ menetapkan $\eta(t) =
\frac{f(t) - f(x) - f'(x)(t-x) - \frac12f''(x)(t-x)^2}
{(t-x)^2}$ untuk $t \neq x$, dan $\eta(x) = 0$: sebab menurut
Taylor--Lagrange $\eta(t) = \frac12\bigl(f''(\xi) - f''(x)\bigr)$ untuk
suatu $\xi$ di antara $t$ dan $x$, jadi $\abs\eta \leq
\norm{f''}_\infty$ dan $\eta(t) \to 0$ ketika $t \to x$ (berkat
kekontinuan $f''$). Setelah uraiannya dijumlahkan terhadap $p_k$ dan
pertanyaan 2--3 dipakai:
\[
n\bigl(B_nf(x) - f(x)\bigr)
= \frac{x(1-x)}{2}f''(x)
+ n\sum_k\eta\Bigl(\frac kn\Bigr)\Bigl(\frac kn -
x\Bigr)^2p_k .
\]
Diberikan $\varepsilon$, pilihlah $\delta$ dengan $\abs\eta \leq
\varepsilon$ pada $\abs{t - x}\leq\delta$. Bagian dekatnya: paling banyak
$\varepsilon\,n\cdot\frac{x(1-x)}n \leq \varepsilon$. Bagian jauhnya:
dengan $C = \norm{f''}_\infty$ dan pertanyaan 23,
\[
n\,C\sum_{\abs{k/n-x}>\delta}\Bigl(\frac kn -
x\Bigr)^2p_k
\leq \frac{nC}{\delta^2}\sum_k\Bigl(\frac kn -
x\Bigr)^4p_k
= \frac{nC}{\delta^2 n^4}\sum_k(k-nx)^4p_k
\leq \frac{C}{\delta^2 n} \longrightarrow 0 .
\]
Karenanya $n(B_nf(x) - f(x)) \to \frac{x(1-x)}2f''(x)$ ---
itulah teorema Voronovskaya. Untuk $f = e_2$ hal ini eksak di setiap
$n$ (\cref{exo:b2:funcseq:7}): jadi langit-langit $\frac1n$ itu nyata.

\textbf{25.} (i) Kepositifan mengubah ketaksamaan \omterm{def:b2:funcseq:def}{titik demi titik}
menjadi ketaksamaan operator: ia memberi batas \omterm{def:b2:nvs:norm}{normanya}, Jensen,
Chebyshev, dan seluruh Korovkin --- sedangkan kelinearan semata tak
membuktikan apa pun di sini. (ii) Kekompakan masuk lewat Heine
(pertanyaan 6 dan 13), lewat keterbatasan $f$, dan lewat kenyataan bahwa
\omterm{def:b2:nvs:norm}{norma} $\norm\cdot_\infty$ itu sendiri hingga. (iii) Tiga fungsi uji
sudah cukup karena kepositifan menyusutkan segalanya menjadi
pengendalian $L_n$ pada satu keluarga $(s-x)^2 = e_2 - 2xe_1 + x^2e_0$,
yang rentangnya sama dengan rentang $e_0, e_1, e_2$. (iv) Bernstein
konvergen dengan laju jujur $\omega(n^{-1/2})$ untuk setiap $f$ yang
\omterm{def:b2:metric:continuity}{kontinu} (dan tajam, pertanyaan 11) tetapi menjenuh pada $\frac1n$ untuk
$f$ yang mulus (pertanyaan 24); sedangkan bab Fourier menjalankan
program yang sama untuk fungsi periodik lewat kernel Fejér --- yakni
operator positif lain dengan keutamaan dan kerendahan hati yang sama.

\end{solution}
